\documentclass[12pt]{article}

\usepackage[spanish,es-nodecimaldot]{babel}

\usepackage{fontspec}

\setmainfont{Arial}

\usepackage[
    top=2.54cm,
    bottom=2.54cm,
    left=2.54cm,
    right=2.54cm
]{geometry}

\usepackage{setspace}
\onehalfspacing

\usepackage{ragged2e}
\justifying

\setlength{\parindent}{0pt}
\setlength{\parskip}{\baselineskip}

\usepackage{titlesec}

\titleformat{\section}
  {\bfseries\fontsize{14}{16}\selectfont}
  {\thesection.}
  {0.5em}
  {}

\titleformat{\subsection}
  {\bfseries\fontsize{14}{16}\selectfont}
  {\thesubsection.}
  {0.5em}
  {}

\titleformat{\subsubsection}
  {\bfseries\fontsize{14}{16}\selectfont}
  {\thesubsubsection.}
  {0.5em}
  {}

\usepackage{fancyhdr}

\pagestyle{fancy}
\fancyhf{}
\fancyfoot[R]{\thepage}
\renewcommand{\headrulewidth}{0pt}

\usepackage{cite}
\usepackage{amsmath,amssymb}
\usepackage{graphicx}
\usepackage{booktabs}
\usepackage{url}
\usepackage{hyperref}

\begin{document}

\begin{center}

{\fontsize{16}{18}\selectfont
\textbf{Análisis del Tiempo de Acceso y Rendimiento de Diferentes Tecnologías de Memoria}
}

\vspace{1cm}

\textbf{Tecnológico de Costa Rica} \\
Escuela de Computación \\
Fundamentos de Organización de Computadoras \\

\vspace{0.5cm}

Víctor Julio Chavarría Ordoñez \\
Luiz Fernando Mendoza Contreras \\
Harold Philippe Umaña Pacheco \\

\vspace{0.5cm}

II Semestre 2026 \\
Fecha: 21 de setiembre de 2026

\end{center}

\vspace{1cm}

\section*{Abstract}
% El abstract debe redactarse en inglés al finalizar la investigación.
This research analyzes memory access time and performance in different memory
technologies, with emphasis on latency and bandwidth. A literature review of
recent scientific studies was conducted, considering technologies such as
DDR4, DDR5, and HBM. The results show that higher bandwidth does not
necessarily imply lower latency, and that memory performance also depends on
the system architecture, memory controller, and application access patterns.
Therefore, there is no single memory technology that provides the best
performance for every workload. Memory selection should consider both latency
and bandwidth according to the requirements of each application.

\textbf{Keywords:}
memory, latency, bandwidth, access time, memory hierarchy, performance

\section{Introducción}
% Responsable(s): por asignar
% Objetivo: presentar el problema, la motivación, el objetivo de la revisión
% y la pregunta de investigación.

% Redacción pendiente.
El rendimiento de los sistemas computacionales actuales no depende solamente
de la capacidad de procesamiento de la CPU, sino también de la rapidez con la
que la jerarquía de memoria puede suministrar los datos requeridos. A pesar de
los avances en la arquitectura de los procesadores, la latencia de acceso a la
memoria DRAM no ha disminuido significativamente en las últimas décadas, por
lo que continúa representando un posible cuello de botella para el rendimiento
del sistema. En arquitecturas multinúcleo, este problema puede intensificarse
debido a la contención producida cuando varios núcleos intentan acceder de
forma simultánea a bancos de memoria compartidos
\cite{alawneh2023dram}.

Una de las principales consecuencias de esta diferencia de velocidades es la 
aparición de \textit{memory stalls}, durante los cuales el procesador debe 
detener temporalmente la ejecución de intrucciones. Debido a esta relación,
la latencia de acceso no solo constituye una característica propia de la 
memoria, sino que también repercute directamente sobre métricas del rendimiento
del procesador. Park et al. muestran que una gestión coordinada entre la CPU 
y los diferentes niveles de la jerarquía de memoria puede reducir los ciclos 
asociados a estas esperas \cite{park2022dvfs}.

El análisis del rendimiento de una memoria tampoco puede limitarse a su ancho 
de banda teórico. En sistemas basados en DDR4 existen retardos asociados con
operaciones internas como los ciclos de refresco, activación y precarga. Estos
factores pueden reducir el throughput que realmente puede aprovechar una aplicación.
Shahbazi y Tsui demostraron experimentalmente esta diferencia al implementar 
cuatro memorias DDR4-2400, donde esto alcanzo un throughput de 135.52~Gbps y 
mejorar aproximadamente un 40\% el ancho de banda \cite{shahbazi2026ddr4}.

Asimismo, el ancho de banda y la latencia deben estudiarse de manera conjunta.
Esmaili-Dokht et al. caracterizan sistemas reales equipados con tecnologías
DDR4, DDR5, HBM2, Y HBM2E mediante curvas de ancho de banda. Sus resultados
muestran que la latencia de acceso varía conforme aumenta la presión sobre el
sistema de memoria. Además, el rendimiento no depende exclusivamente de la
tecnología DRAM, sino también de elementos como la jerarquía de caché, el 
controlador de memoria y la arquitectura del procesador \cite{esmailidokht2024mess}.
A partir de este contexto, el presente trabajo analiza el tiempo de acceso y el
rendimiento de diferentes tecnologías y organizaciones de memoria, con énfasis
en la relación entre latencia, ancho de banda y rendimiento del sistema.


\section{Justificación}
El estudio del tiempo de acceso y del rendimiento de las tecnologías de memoria,
resulta conveniente porque ayuda a comprender factores como la latencia, el 
ancho de banda y la organización de la memoria. Este análisis es importante
ya que ante el crecimiento de aplicaciones que requieren procesar grandes 
cantidades de datos, lo cual depende de una comunicación eficiente entre el 
procesador y la memoria.

Desde una perspectiva más práctica, la investigación permite identificar caracterśticas 
y limitaciones de tecnologías como DDR4, DDR5, y memorias de alto ancho de banda, facilitando
la compresión de los criterios de intervienen en la optimización del sistema de memoria.
Su valor teórico consiste en integrar resultados recientes sobre la relación entre latencia,
ancho de banda y rendimiento. Además, aunque no se pretender desarrollar un nuevo instrumento 
de medición, la revisión y comparación proporciona una metodología organizada para analizar
el comportamiento de distintas tecnologías.

\section{Metodología}
% Responsable(s): por asignar
% Objetivo: describir la revisión sistemática: bases consultadas, cadenas de
% búsqueda, criterios de inclusión/exclusión y proceso de selección de artículos.

% Redacción pendiente.
La investigación se desarrolló mediante una revisión bibliográfica de
publicaciones científicas relacionadas con el tiempo de acceso y el rendimiento
de diferentes tecnologías de memoria. La búsqueda se orientó principalmente a
estudios sobre latencia, ancho de banda, tiempo de acceso, jerarquía de memoria
y tecnologías como DDR4, DDR5 y HBM.

Para la selección de las fuentes se consideraron artículos científicos
publicados entre 2022 y 2026. Se seleccionaron al menos doce artículos
relacionados directamente con el tema, procurando que como mínimo un 60\% de
las publicaciones estuvieran escritas en inglés. Las fuentes fueron localizadas
mediante bases de datos académicas disponibles a través de la Biblioteca del
Tecnológico de Costa Rica, Google académico y IEEEXplore.

Una vez seleccionados los artículos, se identificaron los principales datos
relacionados con la tecnología de memoria, latencia, ancho de
banda, tiempo de acceso, y arquitectura utilizada. Seguidamente, estos datos se 
organizaron y compararon con el propósito de determinar las diferencias de rendimiento 
entre las tecnologías analizadas y de esta manera analizar sus características teóricas y 
su comportamiento en sistemas reales.

Finalmente, los resultados de los diferentes estudios se analizaron de forma
colectiva para analizar limitaciones y avances recientes en los
sistemas de memoria. Este proceso permitió utilizar tanto resultados
cuantitativos como observaciones de los autores para desarrollar el estado del
arte, la comparación de tecnologías y las conclusiones preliminares de la
investigación.


\section{Marco conceptual}

\subsection{Fundamentos}
\input{secciones/04-fundamentos}

\subsection{Tecnologías de memoria}
% Responsable(s): por asignar
% Objetivo: describir las tecnologías de memoria incluidas en la revisión y
% contextualizar sus características relevantes para el análisis de rendimiento.

% Redacción pendiente.
\subsubsection{Memoria DRAM}
\subsubsection{DDR4}
\subsubsection{DDR5}
\subsubsection{Memorias de alto ancho de banda}

\section{Estado del arte}

\subsection{Comparación de rendimiento}
% Responsable(s): por asignar
% Objetivo: comparar cuantitativamente latencia, ancho de banda, tiempo de
% acceso y demás métricas reportadas por los artículos seleccionados.

% Redacción pendiente.
\subsubsection{HBM frente a DDR en la supercomputadora Aurora}

Intel y el Laboratorio Nacional Argonne pusieron a prueba, dentro de la supercomputadora Aurora,
los dos tipos de memoria que conviven en sus nodos de cómputo: la memoria de alto ancho de
banda (HBM) integrada en el procesador Xeon Max, y la DDR tradicional del mismo sistema. El
benchmark que usaron, STREAM Triad, midió el ancho de banda sostenido bajo distintas
combinaciones de clustering (Quad y SNC4) y modos de memoria (Flat y Cache). Con socket
completo en modo SNC4, HBM llegó a 840 GB/s, más del triple de los 245 GB/s que alcanzó
DDR \cite{ibeid2025aurora}. Pero la historia cambia con la latencia: al medirla sin carga usando Intel
Memory Latency Checker, DDR resultó más rápida en todas las configuraciones probadas, con
una ventaja de más del 10\% en modo SNC4. La conclusión de los autores es simple pero útil: ni
HBM ni DDR ganan en todo, así que la mejor opción depende de si la aplicación necesita más
ancho de banda o menos latencia.

\subsubsection{El marco Mess: una vista unificada del sistema de memoria}

Esmaili-Dokht y su equipo ---trabajo que quedó como finalista al mejor artículo en MICRO 2024
--- construyeron algo distinto a lo habitual: en vez de reportar un solo número de latencia o ancho
de banda, propusieron caracterizar la memoria con curvas completas de ancho de banda contra
latencia, a las que llamaron framework Mess. Corrieron su benchmark sobre servidores reales de
Intel, AMD, IBM, Fujitsu y Amazon, además de GPUs NVIDIA, cubriendo DDR4, DDR5, HBM2
y HBM2E. Un dato llamativo: dos servidores con la misma memoria DDR4 pueden tener latencias
muy distintas. Uno con arquitectura Cascade Lake midió 85 ns sin carga, mientras que otro con
Zen2 ---también DDR4, con temporizaciones casi idénticas--- llegó a casi 30 ns más
\cite{esmailidokht2024mess}. Esto deja claro que la latencia no depende solo del chip de memoria, sino de
toda la microarquitectura que lo rodea: caché, interconexión, número de núcleos. Además, algo
útil para el resto del grupo: los autores probaron que simuladores tan usados como gem5 no logran
replicar bien estos comportamientos reales, así que conviene tener cuidado si en algún momento
se citan resultados de simulación en vez de mediciones directas.

\subsubsection{Memoria heterogénea con enlaces de coherencia de caché: CXL, NVLink-C2C e Infinity Fabric}

Con la llegada de CXL, NVLink-C2C e Infinity Fabric, la memoria dejó de limitarse a los canales
DIMM de siempre, y eso es justo lo que investigó este grupo de UC San Diego con su banco de
pruebas Heimdall. Probaron diez tipos de servidores distintos combinando CPUs de varios
fabricantes con DDR5 DIMM, LPDDR5x, HBM3, HBM3e y memoria expandida vía CXL. Los
sistemas DDR5 de cuatro canales saturan su ancho de banda cerca de 107 GiB/s; con ocho canales,
la cifra sube considerablemente antes de que la latencia empiece a dispararse
\cite{wang2024hitchhiker}. Lo más interesante para esta comparación es el costo de usar CXL: un acceso remoto vía este
protocolo puede ser entre 50\% y 300\% más lento que un acceso local a DIMM. En números
concretos, mientras la DRAM local ronda los 100 ns, un dispositivo CXL basado en ASIC se
mueve entre 200 y 300 ns, y los prototipos FPGA llegan hasta 400 ns. Este artículo aporta justo lo
que le faltaba a los otros dos: tecnologías de memoria de última generación ---HBM3, HBM3e,
CXL--- con números concretos de cuánto cuesta, en latencia, expandir la memoria más allá de lo
tradicional.

\subsubsection{Controlador de memoria híbrido con eDRAM y MRAM para Big Data}

El cuarto artículo cambia de enfoque: en vez de medir memoria principal a nivel de sistema
completo, se mete en el diseño del controlador de memoria mismo. Fue presentado en SARTECO
2023, el congreso anual de la Sociedad Española de Arquitectura y Tecnología de Computadores,
por un equipo de la Universitat Politècnica de València. El problema que atacan es conocido: la
memoria no volátil (NVRAM) ofrece mucha más capacidad que la DRAM, pero paga esa
capacidad con latencias de acceso mayores. La solución que proponen es añadir una segunda
caché, más densa que la SRAM habitual, construida con eDRAM o MRAM, que absorbe los fallos
antes de que lleguen a NVRAM. Simulando cargas de trabajo típicas de Big Data, el esquema con
eDRAM bajó la penalización por fallo hasta un 50\%, y el de MRAM hasta un 80\%; a nivel de
sistema completo eso se tradujo en mejoras de rendimiento de 15\% y 23\% respectivamente
\cite{lurbe2023controlador}. Sirve porque es el único de los cuatro que no habla de una tecnología de memoria ya
establecida en el mercado, sino de cómo diseñar la jerarquía de caché para tecnologías que todavía
están emergiendo.

\subsubsection{Comparación cuantitativa}

La siguiente tabla resume los valores numéricos más relevantes reportados en los cuatro artículos,
organizados por tecnología y tipo de medición.

\begin{table}[ht]
\centering
\small
\begin{tabular}{p{4.2cm} p{3cm} p{3cm} p{3.2cm}}
\toprule
\textbf{Tecnología / configuración} &
\textbf{Ancho de banda medido} &
\textbf{Latencia medida} &
\textbf{Fuente} \\
\midrule

HBM (Xeon Max, modo SNC4) &
840 GB/s &
$\sim$121 ns (idle) &
Ibeid et al. (2025) \\

DDR (Xeon Max, modo SNC4) &
245 GB/s &
$\sim$96 ns (idle) &
Ibeid et al. (2025) \\

DDR4 (servidor Cascade Lake) &
--- &
85 ns (sin carga) &
Esmaili-Dokht et al. (2024) \\

DDR5 (Graviton 3) &
--- &
129 ns (sin carga) &
Esmaili-Dokht et al. (2024) \\

HBM2 (A64FX) &
--- &
122 ns (sin carga) &
Esmaili-Dokht et al. (2024) \\

DDR5 DIMM (4 canales, punto de saturación) &
107 GiB/s &
--- &
Wang et al. (2024) \\

Memoria CXL remota (ASIC) vs. DIMM local &
--- &
200--300 ns vs. $\sim$100 ns &
Wang et al. (2024) \\

Caché híbrida eDRAM/MRAM vs. SRAM sola &
--- &
$-50\%$ a $-80\%$ en penalización por fallo &
Lurbe et al. (2023) \\

\bottomrule
\end{tabular}
\caption{Comparación cuantitativa de tecnologías y configuraciones de memoria.}
\label{tab:comparacion-memoria}
\end{table}

Viendo estos números juntos, se repite una misma idea en los cuatro artículos: no existe una
memoria que gane en ancho de banda y en latencia al mismo tiempo. HBM y sus variantes se
llevan la ventaja en ancho de banda ---algo clave para cargas de trabajo muy paralelas, como HPC
o entrenamiento de modelos de IA--- pero pierden entre 10\% y 20\% de latencia frente a DDR
\cite{ibeid2025aurora,esmailidokht2024mess}. Salir de los canales DIMM tradicionales también
tiene un precio, sea por CXL o por una jerarquía de caché adicional: CXL suma entre 100 ns y 300
ns sobre un acceso local \cite{wang2024hitchhiker}, mientras que las cachés híbridas con eDRAM o
MRAM hacen justo lo contrario, intentan esconder la latencia de la NVRAM a cambio de un
controlador de memoria más complejo \cite{lurbe2023controlador}. Y algo que vale la pena recordar: la
latencia real que se mide en un sistema depende tanto de la memoria como de todo lo que la rodea
---caché, interconexión, cantidad de núcleos---, razón por la cual dos servidores con la misma
DDR4 pueden diferir en casi 30 ns \cite{esmailidokht2024mess}.

\subsection{Investigaciones recientes}
% Responsable(s): por asignar
% Objetivo: sintetizar tendencias, enfoques actuales y hallazgos recurrentes
% identificados en los artículos seleccionados.

% Redacción pendiente.
En conjunto, estos cuatro artículos cubren tanto memoria principal ya consolidada ---DDR4,
DDR5, HBM, HBM2, HBM3--- como tecnologías más nuevas de expansión y caché de memoria
---CXL, eDRAM, MRAM---, todos con datos medidos en hardware o en simulación detallada, y
todos publicados entre 2023 y 2025. Esa tensión entre ancho de banda y latencia que aparece una
y otra vez, junto con el costo que implica salir de los canales DIMM convencionales, es la
evidencia que este bloque aporta al análisis comparativo del grupo.
De manera complementaria, Eyerman et al. proponen analizar el comportamiento
de la memoria DRAM mediante \textit{bandwidth stacks} y \textit{latency stacks},
con el objetivo de identificar las causas que limitan el aprovechamiento del
ancho de banda y aumentan la latencia. Su trabajo muestra que el rendimiento
real de DRAM está condicionado por restricciones temporales, paralelismo interno
y esperas en el controlador de memoria, por lo que las especificaciones teóricas
no representan necesariamente el comportamiento observado en una aplicación
\cite{eyerman2022dram}.

\subsection{Análisis y discusión}
% Responsable(s): por asignar
% Objetivo: interpretar los resultados, contrastar ventajas y limitaciones,
% discutir diferencias entre estudios y relacionar los hallazgos con IC-1400.

% Redacción pendiente.
Los resultados analizados muestran que el rendimiento de una memoria no puede
evaluarse utilizando únicamente una característica. Existe un compromiso entre
latencia y ancho de banda, ya que tecnologías como HBM pueden proporcionar un
ancho de banda considerablemente mayor que DDR, pero no necesariamente una
menor latencia \cite{ibeid2025aurora}. Por esta razón, la tecnología más
adecuada depende de las necesidades y del patrón de acceso de cada aplicación.

Además, el rendimiento observado depende tanto de la memoria como de la
arquitectura del sistema. Esmaili-Dokht et al. muestran que sistemas que
utilizan una misma tecnología de memoria pueden presentar diferencias
importantes de latencia debido a factores como la jerarquía de caché, la
interconexión y la microarquitectura \cite{esmailidokht2024mess}. De manera
similar, Velten et al. muestran que la organización de la jerarquía de memoria,
la ubicación de los datos y las características de acceso local o remoto pueden
producir diferencias importantes en latencia y ancho de banda en sistemas de
servidores \cite{velten2022memory}. Las restricciones internas y el controlador
de memoria también pueden limitar el aprovechamiento del ancho de banda
disponible \cite{eyerman2022dram}.

En conjunto, los artículos revisados indican que no existe una tecnología de
memoria que presente el mejor rendimiento en todas las situaciones. Algunas
aplicaciones pueden beneficiarse principalmente de un mayor ancho de banda,
mientras que otras dependen más de reducir los tiempos de acceso. Por ello, la
evaluación de una tecnología de memoria debe considerar conjuntamente la
latencia, el ancho de banda y las características de la arquitectura donde será
utilizada.

Al comparar estas tecnologías, también se observa que la evolución
de DDR4 y DDR5 está orientada principalmente a aumentar la capacidad de 
transferencia. DDR5 puede ofrecer mayor ancho de banda gracias a cambios 
en su organización y a velocidades de transferencia más elevadas 
\cite{steiner2022dramsys,choi2024ddr5}. Sin embargo, esto no significa
que DDR4 deje de ser competitiva, ya que determinadas cargas sensibles a 
la latencia pueden obtener resultados favorables con esta tecnología 
\cite{steiner2022dramsys}. Trabajos como el de Shahbazi y Tsui muestran que 
una mejor organización del controlador puede mejorar notablemente el 
rendimiento de DDR4, incluso sin cambiar la memoria en sí \cite{shahbazi2026ddr4}. 
Esto confirma que el rendimiento final depende tanto de la generación de memoria 
como de su integración y gestión en el sistema.

\section{Conclusiones preliminares}
% Responsable(s): por asignar
% Objetivo: responder la pregunta de investigación, resumir los hallazgos
% principales y señalar limitaciones o líneas de trabajo futuras.

% Redacción pendiente.
Los estudios revisados permiten concluir de manera preliminar que el
rendimiento de un sistema de memoria no depende únicamente de la tecnología
utilizada, sino también de factores como la latencia, el ancho de banda, la
organización de la jerarquía de memoria y la arquitectura del sistema. Los
resultados analizados muestran que tecnologías con un mayor ancho de banda,
como HBM, no necesariamente presentan los menores tiempos de acceso, por lo
que no existe una única alternativa que sea superior para todas las cargas de
trabajo \cite{ibeid2025aurora,esmailidokht2024mess}.

Asimismo, las investigaciones recientes muestran que el rendimiento real puede
diferir considerablemente de los valores teóricos debido al controlador de
memoria, la microarquitectura, la ubicación de los datos y los patrones de
acceso de las aplicaciones \cite{eyerman2022dram,velten2022memory}. Este
comportamiento también se observa en tecnologías de alto ancho de banda, cuyo
rendimiento efectivo depende de la manera en que se organizan y realizan los
accesos a memoria \cite{huang2022shuhai}. Tecnologías como CXL y las memorias
heterogéneas amplían las posibilidades de capacidad y organización del sistema,
aunque también pueden introducir costos adicionales de latencia
\cite{wang2024hitchhiker,lurbe2023controlador}.

A partir de la evidencia recopilada, la investigación resulta viable, ya que
existen estudios recientes con mediciones cuantitativas que permiten comparar
diferentes tecnologías y configuraciones de memoria. De forma preliminar, puede
establecerse que la selección de una tecnología de memoria debe realizarse
considerando conjuntamente el tiempo de acceso, la latencia, el ancho de banda
y las necesidades específicas de cada aplicación. Estos resultados serán
complementados con el análisis de las demás tecnologías incluidas en la
investigación.

\renewcommand{\refname}{Referencias bibliográficas}

\bibliographystyle{IEEEtran}
\bibliography{referencias}

\end{document}